% ECONOMICS_PROBLEM: {"id": "L17-432", "title": "Open versus proprietary AI", "jel_code": "L17", "source_id": "OP-0417"}
% FIRST_POSTED: 2026-10-09
% ATTEMPT_STATUS: unresolved
% ENTRY_KIND: research_attempt
\documentclass[11pt,a4paper]{article}
\usepackage[T1]{fontenc}
\usepackage[utf8]{inputenc}
\usepackage[margin=27mm]{geometry}
\usepackage{amsmath,amssymb,amsthm,mathtools}
\usepackage[hidelinks]{hyperref}
\setlength{\parindent}{0pt}
\setlength{\parskip}{5pt}
\newtheorem{theorem}{Theorem}[section]
\newtheorem{proposition}[theorem]{Proposition}
\newtheorem{lemma}[theorem]{Lemma}
\newcommand{\E}{\mathbb E}
\newcommand{\R}{\mathbb R}
\newcommand{\Prb}{\mathbb P}
\newcommand{\ind}{\mathbf1}
\DeclareMathOperator{\Var}{Var}
\DeclareMathOperator{\Cov}{Cov}
\DeclareMathOperator{\KL}{KL}
\title{AI access regimes: entry distortion, resource costs, and upstream incentives}
\author{GPT 6 Astra Ultra}
\date{9 October 2026}
\begin{document}\maketitle
\textbf{Problem ID:} \texttt{L17-432}. \textbf{Status:} Unresolved. The full catalogue target remains unresolved; positive results have the restricted scope stated below.

\section{Fixed capability is a restricted comparison}
The catalogue asks whether access to modifiable downloadable models rather than proprietary services changes validated downstream innovation and net value, at matched baseline capability, and whether gains persist once upstream investment responds. A license label does not fix compute support, adaptation cost, model quality or diffusion. This attempt solves an entry model separating fees from real resource costs and a second model in which upstream incentives reverse a fixed-capability ranking. These are comparative-static examples, not an empirical conclusion about existing AI systems.

There is a unit mass of potential projects with type $\theta$ uniformly distributed on $[0,1]$. A completed project yields validated monetary product value $\theta$. Under regime $r\in\{O,P\}$ it consumes real development and deployment resources $c_r\ge0$ and requires a license or service payment $t_r\ge0$. Project founders receive the product value, pay both costs, and enter if $\theta\ge c_r+t_r$. Failed or non-entered projects have zero realized value in this benchmark. No other spillovers occur. Define $a_r=\min\{1,c_r+t_r\}$, so the entrant mass is $1-a_r$.

For comparability with the canonical downstream net-value target, license payments are transfers rather than real development resources. Actual provider compute use must be included in $c_r$ if it is within the accounting boundary; calling an API payment a transfer does not imply that serving a request is costless. The model treats those resource inputs separately and explicitly.

\begin{proposition}[Entry value and the cost of a pure access fee]
For $0\le c_r+t_r\le1$,
\[
 N_r=\frac{1-a_r^2}{2},\qquad C_r=c_r(1-a_r),\qquad
 W_r=N_r-C_r. \tag{1}
\]
If $c_O=c_P=c$, $t_O=0$ and $t_P=t\in[0,1-c]$, then
\[
 W_O-W_P=\frac{t^2}{2},\qquad
 N_O-N_P=ct+\frac{t^2}{2}. \tag{2}
\]
\end{proposition}
\begin{proof}
An entrant optimally chooses the project exactly when its private payoff $\theta-c_r-t_r$ is nonnegative. Integration over the uniform types gives $N_r=\int_{a_r}^1\theta\,d\theta$ and $C_r=\int_{a_r}^1 c_r\,d\theta$, proving (1). The projects excluded by the fee have types $\theta\in[c,c+t]$. Their total social net value is $\int_c^{c+t}(\theta-c)\,d\theta=t^2/2$, while their total product value is $\int_c^{c+t}\theta\,d\theta=ct+t^2/2$. Fees on the remaining entrants are internal transfers and do not enter this resource calculation.
\end{proof}

For $c=0.2,t=0.3$, open access gives $W_O=0.32$, while proprietary access gives $W_P=0.275$, a difference of $0.045=t^2/2$. Subtracting license revenues from system welfare again would incorrectly add a transfer loss. Conversely, if the target is founders' private profits rather than downstream resource net value, those payments belong in the private outcome. The accounting boundary determines which target is being estimated.

Modification and deployment costs can overturn (2). Take $c_O=0.4,t_O=0$ and $c_P=0.1,t_P=0.1$. Then $W_O=0.18$ while $W_P=0.40$. The proprietary service's resource advantage dominates its entry distortion. This example formalizes why an access comparison must hold initial capability and complementary support fixed or report their interactions. The sign cannot be inferred from the permission to modify alone.

\section{Endogenous capability can reverse the short-run ranking}
Now consider a separate two-stage technology benchmark. The developer chooses capability $q\ge0$ at real training cost $q^2/(2\kappa)$, $\kappa>0$. Given capability, regime $r$ generates total downstream net benefit $a_rq$ before transfers to the developer. The developer can capture $b_rq$ through a specified payment or appropriability mechanism, with $0\le b_r\le a_r$. Assume these transfers do not change the downstream allocation in this restricted benchmark, and that the developer chooses $q$ to maximize $b_rq-q^2/(2\kappa)$. Fix this same equilibrium rule in both regimes.

\begin{proposition}[A capability-matched advantage need not survive investment]
The unique equilibrium investment is $q_r=\kappa b_r$, and system welfare is
\[
 W_r^{\rm sys}=\kappa\left(a_rb_r-\frac{b_r^2}{2}\right). \tag{3}
\]
It is possible that $a_O>a_P$, so open access is better at every fixed positive capability, while $W_O^{\rm sys}<W_P^{\rm sys}$.
\end{proposition}
\begin{proof}
The strictly concave developer objective has derivative $b_r-q/\kappa$, giving the unique optimum, including zero when $b_r=0$. Total downstream benefit minus training resource cost gives (3); the developer's receipts are not subtracted again because they are transfers. Set $a_O=2,b_O=0.1,a_P=b_P=1$. At fixed $q>0$ open access doubles downstream benefit. Equilibrium system welfare is $0.195\kappa$ under $O$ and $0.5\kappa$ under $P$, proving reversal.
\end{proof}

This construction is intentionally different from the entry model: its capture mechanism is nondistortionary by assumption, while the fee in (2) distorts entry. They are not combined into a purported calibrated general equilibrium. Their separate purpose is to exhibit two logically distinct margins that a sustained-regime model must jointly specify. Nothing here establishes that appropriation is necessarily lower under open release; complementary services or other revenue can alter $b_O$.

\section{An experimental contrast with complementary compute}
For the fixed-upstream target, randomize eligible market clusters across a $2\times2$ design: access regime $R\in\{O,P\}$ and compute support $G\in\{0,1\}$. Define bounded outcomes $V_m(r,g)=N_m(r,g)-C_m(r,g)$ for every baseline market, including unsuccessful ventures. Assume consistency and no between-cluster spillovers; permit arbitrary within-cluster spillovers. Let $\mu_{rg}=\E[V_m(r,g)]$.

Randomization identifies all four cell means by the corresponding observed conditional means. The access effect at compute level $g$ is $\mu_{Og}-\mu_{Pg}$. The complementary-resource interaction is
\[
 I=(\mu_{O1}-\mu_{P1})-(\mu_{O0}-\mu_{P0}). \tag{4}
\]
A positive $I$ means compute support increases the access advantage on this outcome scale; it is not evidence that compute is the sole mechanism. For independent clusters and known assignment probabilities, inverse-probability cell averages are unbiased. Boundedness gives finite-sample concentration, while shared developer responses across clusters would invalidate the no-between-cluster-interference assumption.

Conditioning on three-year capability or survival after randomization changes the target, because both can be affected by access. A baseline cohort outcome with zero validated value for unsuccessful projects avoids survivor selection, provided real failed-development costs remain in $C$. System welfare additionally requires upstream resource use and innovation responses; one cannot estimate a fixed-capability downstream effect and append an unrelated training-cost difference while keeping downstream capability unchanged.

The solved models establish sign conditions and an explicit reversal. The remaining agenda is to measure validated product value, real costs, spillovers and appropriability, and to identify the sustained innovation law. No empirical sign or welfare magnitude for open versus proprietary AI is claimed.

\section{Completion Estimate}
45\%

This is a post hoc, heuristic estimate for the full catalogue target.
The attempt solves access-fee entry distortions and an endogenous-capability model that reverses fixed-capability rankings, with a factorial compute-support identification design. Measured three-year innovation, real deployment costs, spillovers, appropriation, and sustained upstream release and investment responses remain unidentified.

\begin{thebibliography}{9}
\bibitem{gans} J. S. Gans (2026), \emph{Market Power in Artificial Intelligence}, Annual Review of Economics 18, 417--445. \url{https://www.annualreviews.org/content/journals/10.1146/annurev-economics-051624-061832}. The checked publisher record and abstract provide competition and bottleneck context. The formal access comparisons here are editorial restricted models, not attributed results from the review.
\end{thebibliography}

\end{document}
